% Fragmento público generado desde propietarios íntegros.
% Residencia: 52.11.1.
% Fuente: ../incoming/radion_monografia_20260827/manuscrito/sections/12_aritmetica_borde.tex:1-159
\noindent La frontera nonádica conserva el defecto suma--producto, la vacancia y el acarreo de completación.\par\medskip

\subsubsection{Radical \(3\! -\! 6\! -\! 9\) del borde}

\paragraph{Generación y nilpotencia del radical nonádico}

En el anillo \(\Z/9\Z\), el ideal
\begin{equation}
\mathfrak m=(3)=\{0,3,6\}
\label{eq:radical-369}
\end{equation}
satisface
\begin{equation}
\mathfrak m^2=(9)=0.
\label{eq:m-square-zero}
\end{equation}
Es el sector nilpotente del borde nonádico. La órbita visible de sus tres
marcas se publica como
\[
369\longrightarrow693\longrightarrow936\longrightarrow369.
\]
El retorno no implica que todos los estados sean iguales: la permutación
visible cierra mientras los cociclos de ruta pueden avanzar.

La doble convención
\[
9\equiv0\pmod9,
\qquad
\operatorname{dr}_9(9)=9
\]
expresa dos lectores: cero residual algebraico y cierre visible de raíz
digital. Por eso un nueve puede absorber torsión de borde sin convertirse en
un cero específico de la función zeta.

\paragraph{Defecto suma--producto \(459-405=54\)}

Las hojas APP producen totales complementarios
\[
S_+=405,\qquad S_\times=459,\qquad S_\times-S_+=54.
\]
El \(54\) es un defecto exacto de las tablas, y reaparece como exponente de
contracción por la identidad \(9\cdot6=54\). Sin embargo, cada aparición
requiere su mapa. La distancia focal \(d_1=0.544545\ldots\) no es igual a
\(0.54\), y el análisis de invariantes no ha producido una flecha que la
identifique con el defecto \(54\).

La potencia
\[
9^9=387,420,489
\]
genera una escala donde
\begin{equation}
\operatorname{ord}_{10}\bigl(9^{9^9}\bigr)=369,693,099,
\qquad
\#\,\text{dígitos}=369,693,100.
\label{eq:369693}
\end{equation}
El bloque \(369|693\) aparece en el par orden--longitud; no se afirma que sea
el prefijo decimal de la potencia.

\subsubsection{Teoría aritmética de las vacancias decimales}

\paragraph{Pendiente asintótica de vacancia}

La separación entre la década y la base nonádica es
\begin{equation}
\delta_{\mathrm{vac}}=\log_{10}\frac{10}{9}.
\label{eq:delta-vac}
\end{equation}
El contador acumulado
\begin{equation}
V_9(n)=\left\lceil n\delta_{\mathrm{vac}}\right\rceil
\label{eq:vacancy-count}
\end{equation}
registra cuántas posiciones de corrección necesita la publicación decimal de
una evolución nonádica. Sus saltos forman una palabra binaria
\[
z_n=V_9(n)-V_9(n-1)\in\{0,1\}.
\]
Los huecos se distribuyen con longitudes \(21/22\), los retornos con
longitudes \(6/7\), y el determinante combinado es \(15\). Estos números
proceden del lector de vacancia, no del núcleo Barbero aunque el \(15\)
coincida con \(\binom62\).

\paragraph{Polarización y corriente aritmética de la red nonádica}

Definimos
\begin{equation}
P_N=\sum_{n=1}^N(z_n-\delta_{\mathrm{vac}}),
\qquad
J_N=P_N-P_{N-1}=z_N-\delta_{\mathrm{vac}}.
\label{eq:arith-current}
\end{equation}
La secuencia equilibrada satisface \(|P_N|<1\). La \enquote{corriente}
\(J_N\) es aquí una corriente aritmética: mide creación y absorción de
vacancias respecto a la pendiente media. Sólo un transductor dimensional
posterior puede publicarla como densidad de corriente electromagnética.

La intuición física es potente porque la misma forma algebraica separa un
promedio par y una fluctuación orientada. El rigor exige conservar el tipo:
\[
J_N\in\R\quad\text{adimensional},
\qquad
j^\mu\in\text{recta física de corriente}.
\]

\subsubsection{Completación \(729\to1000\) y conservación del acarreo}

\paragraph{Descomposición \(1000=729+243+27+1\)}

La completación
\[
1000=729+243+27+1
\]
conserva la unidad mediante cuatro escalas ternarias. La corona total es
\[
1000-729=271,
\]
mientras el último entero antes del cierre satisface
\[
999=729+270.
\]
La diferencia entre \(270\) y \(271\) distingue frontera ocupada de unidad
completante. El mismo cuidado evita identificar una vacancia con un cero o un
acarreo con un residuo real.

\paragraph{Residencia del acarreo del radión en la corona decimal}

Los operandos del radión comienzan, en bloques de tres cifras, por
\[
\Phi_T:\quad000|027|455|906|837|565|\cdots,
\]
\[
D_E:\quad000|027|455|010|034|743|\cdots.
\]
La resta APP da
\[
\epsone:\quad000|000|000|896|802|822|\cdots.
\]
El pequeño tamaño del resultado no se obtiene truncando los sumandos bajos.
Se obtiene porque dos publicaciones distintas comparten un largo prefijo y el
operador \(\operatorname{SubAcarreo}_{1000}\) conserva los préstamos necesarios
para restarlas exactamente.

La palabra \emph{memoria} tiene aquí dos niveles:
\begin{enumerate}
\item \(q_{\mathrm{mem}}\) registra la profundidad de retornos nonádicos;
\item \(c_i\) registra los préstamos locales de la resta decimal.
\end{enumerate}
El primero organiza la hélice; el segundo publica el decimal. \(\epsR\) usa
ambos porque sus operandos proceden de cilindros a profundidad y su resta se
realiza con acarreo, pero no es igual a ninguno de los contadores.

\begin{thesisbox}
El borde no es un lugar donde la teoría pierde información. Es el lugar donde
la información cambia de representación: fase que vuelve, memoria que
asciende, cilindro que se contrae, vacancia que se registra y acarreo que
transporta la diferencia.
\end{thesisbox}
